% SOURCE_AUTHORITY: sga5_fr_workpass.tex, lines 12927--12941 (LF-normalized unit).
% SOURCE_SHA256: 791F4EFFC5E02832D5D77ED03518C8156D6F07E4C8238B03545DB93D883FBB28
\subsection*{7. Fórmula de Euler--Poincaré}

\begin{statement}{Teorema 7.1}
Sejam $C$ uma curva algébrica conexa, lisa e própria sobre o corpo algebricamente fechado $k$, $\eta$ o ponto genérico de $C$, $\Delta\subset D^+(A)$ uma subcategoria triangulada de $D^+(A)$ estável sob somandos diretos, e $F$ um complexo de feixes de $A$-módulos tal que sejam satisfeitas as condições i), ii) de 7.1. Nestas condições, tem-se $\RG_C(F)\in\Ob(\Delta)$, e é válida a fórmula seguinte (fórmula de Euler--Poincaré):
\[
\clK{\Delta}(\RG_C(F))
=
\EP(C)\clK{\Delta}(F_{\bar\eta})
-\sum_{x\in C(k)}\varepsilon_x^\Delta(F),
\tag{7.2}
\]
onde os $\varepsilon_x^\Delta(F)$ são definidos em 6.2 (fórmula 6.2.2).
\end{statement}

A demonstração deste teorema será feita em várias etapas.
